\subsection{State Processing and Closed-Loop Timing}\label{sec:3.3}

Opening, do not touch it yet.

\subsubsection{Vicon State Processing and Filtering}\label{sec:3.3.1}

The Vicon system supplies the pose of the tracked airframe rigid body: three
positions and three attitude angles, \((x,y,z,\phi,\theta,\psi)\). The
marker-based rigid-body reconstruction itself was introduced in Section~3.1, so
this subsection only defines how that 6-DoF measurement is converted into the
state used by the controller.

The Vicon stream was configured with \(+X\) forward, \(+Y\) left, and \(+Z\)
up, matching the arena convention used by the flight tests. A short frame
calibration was performed before the glider calibration and closed-loop runs:
the airframe was held at a known arena point, valid Vicon samples were averaged,
and the resulting offset was written to the active calibration record. If
\(\mathbf{p}_{V,k}\) is the raw Vicon position at sample \(k\),
\(\bar{\mathbf{p}}_V\) is the mean raw position during this calibration, and
\(\mathbf{p}^{\star}\) is the known arena point, the online position transform
is
\begin{equation}
    \mathbf{p}_{w,k}
    =
    R_z(\psi_0)\mathbf{p}_{V,k}+\mathbf{b},
    \qquad
    \mathbf{b}
    =
    \mathbf{p}^{\star}-R_z(\psi_0)\bar{\mathbf{p}}_V ,
    \label{eq:vicon_position_transform}
\end{equation}
where \(R_z(\psi_0)\) is an optional yaw-alignment rotation and \(\mathbf{b}\)
is the calibrated position offset. The same calibration record is saved with a
hash in the flight logs, so that each throw can be associated with the transform
used during data collection. The orientation check then fixed the attitude signs
used by the controller,
\begin{equation}
    \begin{bmatrix}\phi_k\\\theta_k\\\psi_k\end{bmatrix}
    =
    \operatorname{wrap}_{\pi}
    \left(
    \begin{bmatrix}
    1 & 0 & 0\\
    0 & -1 & 0\\
    0 & 0 & -1
    \end{bmatrix}
    \begin{bmatrix}\phi_{V,k}\\\theta_{V,k}\\\psi_{V,k}\end{bmatrix}
    +
    \begin{bmatrix}0\\0\\\psi_0\end{bmatrix}
    \right),
    \label{eq:vicon_attitude_transform}
\end{equation}
so roll is unchanged, while nose-up pitch and nose-right yaw are positive in the
controller convention.

The controller state contains position, attitude, body velocity, body angular
rate, and estimated actuator state,
\begin{equation}
    \mathbf{x}_k =
    \begin{bmatrix}
    x_w & y_w & z_w &
    \phi & \theta & \psi &
    u & v & w &
    p & q & r &
    \delta_a & \delta_e & \delta_r
    \end{bmatrix}^{T}.
    \label{eq:controller_state_from_vicon}
\end{equation}
Only the first six components are measured directly by Vicon. The remaining
components are derived online from the filtered pose history and the commanded
surface history.

For translational velocity, the sample interval \(\Delta t_k\) is taken from
the Vicon frame number and frame rate when available, with host time used only
as a fallback. The raw finite-difference velocity is
\begin{equation}
    \tilde{\mathbf{v}}_{w,k}
    =
    \frac{\mathbf{p}_{w,k}-\mathbf{p}_{w,k-1}}{\Delta t_k},
    \label{eq:vicon_raw_velocity}
\end{equation}
and the online state estimator applies a one-pole low-pass update,
\begin{equation}
    \mathbf{v}_{w,k}
    =
    \mathbf{v}_{w,k-1}
    +
    \alpha_k
    \left(
    \tilde{\mathbf{v}}_{w,k}-\mathbf{v}_{w,k-1}
    \right),
    \qquad
    \alpha_k = 1-\exp\left(-2\pi f_c\Delta t_k\right),
    \label{eq:vicon_velocity_filter}
\end{equation}
with \(f_c=8~\mathrm{Hz}\). This cutoff is the single online state-filter
choice used in the flight-test runtime and in the latency model: it suppresses
frame-to-frame differentiation noise while adding a bounded feedback delay of
approximately \(20~\mathrm{ms}\). The filtered world velocity is then expressed
in body coordinates as
\begin{equation}
    \begin{bmatrix}u_k\\v_k\\w_k\end{bmatrix}
    =
    C_{wb,k}^{T}
    \operatorname{diag}(1,1,-1)
    \mathbf{v}_{w,k},
    \label{eq:vicon_body_velocity}
\end{equation}
where \(C_{wb,k}\) is the body-to-world direction cosine matrix formed from the
corrected attitude in Eq.~\eqref{eq:vicon_attitude_transform}.

Angular rates are estimated from rotation increments on \(SO(3)\), rather than
by differencing Euler angles. This keeps the rate estimate consistent with the
geometry of rotations, following the same motivation as geometric
Savitzky--Golay filtering on \(SO(3)\)~\cite{jongeneel2022geometric}. For two
successive corrected attitudes,
\begin{equation}
    Q_k = C_{wb,k-1}^{T}C_{wb,k},
    \qquad
    \tilde{\boldsymbol{\omega}}_{b,k}
    =
    \frac{1}{\Delta t_k}\log(Q_k)^{\vee}.
    \label{eq:vicon_so3_rate}
\end{equation}
The online implementation then clips implausible one-frame rates, applies a
seven-frame robust window to downweight isolated rotation spikes, and filters
the candidate body-rate vector with the same \(8~\mathrm{Hz}\) one-pole update.
The launch gate requires the observer confidence to be at least 0.65 for two
consecutive approved frames, so a single noisy Vicon frame cannot start active
control.

Finally, the three control-surface states are not measured by rigid-body Vicon
tracking. They are propagated from the command history using the first-order
actuator model used later for the timing analysis,
\begin{equation}
    \boldsymbol{\delta}_k
    =
    \boldsymbol{\delta}_{k-1}
    +
    \left(
    1-\exp\left(-\Delta t_k/\boldsymbol{\tau}_{\delta}\right)
    \right)
    \left(
    \boldsymbol{\delta}^{cmd}_k-\boldsymbol{\delta}_{k-1}
    \right).
    \label{eq:vicon_actuator_state}
\end{equation}
Thus the Vicon subsystem supplies a calibrated pose stream, while the controller
receives the full 15-state vector in Eq.~\eqref{eq:controller_state_from_vicon}
with velocity, body-rate, and actuator components estimated explicitly.

\subsubsection{Control Surface Response and Implementation Timing}\label{sec:3.3.2}

\vspace{-8pt}

\noindent
\begin{minipage}[t]{0.52\textwidth}
\vspace{0pt}
\vspace{3.9pt}

\hspace{1.5em}The results show that the command-to-surface path has a repeatable effective delay. Across the accepted tests, the median response onset, defined by the 10\% crossing, was approximately \(73~\mathrm{ms}\). The nominal effective latency, defined by the 50\% crossing, was approximately \(108~\mathrm{ms}\). The near-final response, defined by the 90\% crossing, was approximately \(130~\mathrm{ms}\). A conservative upper bound based on the largest valid observed \(t_{50}\) was approximately \(151~\mathrm{ms}\).
\end{minipage}
\hfill
\begin{minipage}[t]{0.44\textwidth}
\vspace{0pt}
\captionsetup{type=figure}
\captionsetup[subfigure]{justification=centering,aboveskip=0pt,belowskip=0pt}
\centering

\begin{subfigure}[t]{\linewidth}
    \centering
    \includegraphics[width=\linewidth,keepaspectratio]{Figures/3/3/Control Surface Mode.jpg}
    \vspace{-0.27cm}
    \caption{Individual control surface tracking.}
    \label{fig:25.1.1.1}
\end{subfigure}

\vspace{0.20cm}

\begin{subfigure}[t]{\linewidth}
    \centering
    \includegraphics[width=\linewidth,keepaspectratio]{Figures/3/3/Rigid Body Mode.jpg}
    \vspace{-0.27cm}
    \caption{Airframe rigid-body tracking.}
    \label{fig:25.1.1.2}
\end{subfigure}

\vspace{-0.15cm}
\caption{Vicon tracking configurations.}
\label{fig:25.1.1}
\end{minipage}

\vspace{2.0pt}
